\item \textbf{Carga e inspección de un catálogo estelar}.

\begin{tcolorbox}[colback=blue!2!white,colframe=blue!55!black,title={Enunciado},
                  before skip=3pt,after skip=5pt,boxsep=3pt,top=3pt,bottom=3pt]
\small
\raggedright
Realice las siguientes operaciones:

\begin{itemize}
\setlength{\topsep}{4pt}
\setlength{\itemsep}{3pt}
\setlength{\parsep}{0pt}
\item[(a)] Importe \texttt{pandas} y cargue los datos en un \texttt{DataFrame}
           llamado \texttt{estrellas} mediante \texttt{pd.read\_csv()} y la
           siguiente ruta:
           \begin{center}
           \url{https://raw.githubusercontent.com/jperaltac/pcfi161/main/tareas/tarea03/estrellas.csv}
           \end{center}
\item[(b)] Muestre las primeras cinco filas del \texttt{DataFrame}.
\item[(c)] Determine el número de filas y columnas.
\item[(d)] Muestre los nombres de las columnas.
\item[(e)] Seleccione solamente las columnas \texttt{estrella}, \texttt{tipo} y
           \texttt{temperatura\_K}.
\item[(f)] Calcule la temperatura media para cada tipo estelar.
\item[(g)] Determine la magnitud mínima registrada para cada tipo estelar.
\item[(h)] Cuente cuántas estrellas fueron observadas con cada instrumento.
\end{itemize}
\end{tcolorbox}

\begin{solution}
\begin{pythonjupyter}
# (a)
import pandas as pd

url = "https://raw.githubusercontent.com/jperaltac/pcfi161/main/tareas/tarea03/estrellas.csv"
estrellas = pd.read_csv(url)

# (b)
print(estrellas.head())

# (c)
print(estrellas.shape)

# (d)
print(estrellas.columns)

# (e)
print(estrellas[["estrella", "tipo", "temperatura_K"]])

# (f)
temperatura_media = estrellas.groupby("tipo")["temperatura_K"].mean()
print(temperatura_media)

# (g)
magnitud_minima = estrellas.groupby("tipo")["magnitud"].min()
print(magnitud_minima)

# (h)
conteo_instrumentos = estrellas["instrumento"].value_counts()
print(conteo_instrumentos)
\end{pythonjupyter}
\end{solution}
